\documentclass[14pt,aspectratio=169,xcolor=dvipsnames]{beamer}
\usetheme{SimplePlus}
\usepackage{booktabs}
\usepackage{minted}
\usepackage{mathtools}

\title[short title]{Clase 6: Polimorfismo y desarrollo de código}
\subtitle{}
\author[NA Barnafi] {Nicolás Alejandro Barnafi Wittwer}
\institute[UC|CMM] 
{
    Pontificia Universidad Católica de Chile \\
    Centro de Modelamiento Matemático
}

\titlegraphic{
    \vspace{-2.4cm}
    \begin{flushright}
      \includegraphics[height=2.5cm]{../images/logos/puc.png} 
    \end{flushright}
}
\date{}

\begin{document}
%%%%%%%%%%%%%%%%%%%%%%%%%%%%%%%%%%%%%%%%%%%%%%%%%%%%%%%
\begin{frame}
    \maketitle
\end{frame}
%%%%%%%%%%%%%%%%%%%%%%%%%%%%%%%%%%%%%%%%%%%%%%%%%%%%%%%
%%%%%%%%%%%%%%%%%%%%%%%%%%%%%%%%%%%%%%%%%%%%%%%%%%%%%%%
\section{Clases y objetos}
%%%%%%%%%%%%%%%%%%%%%%%%%%%%%%%%%%%%%%%%%%%%%%%%%%%%%%%
%%%%%%%%%%%%%%%%%%%%%%%%%%%%%%%%%%%%%%%%%%%%%%%%%%%%%%%
\begin{frame}[fragile]\frametitle{Motivación}
    \begin{minted}{python}
>>> df = pd.DataFrame() # Constructor
>>> type(df)
  <class 'pandas.core.frame.DataFrame'>
    \end{minted}
    \begin{itemize}
        \item Crear nuestros propios tipos de datos
        \item Crear abstracciones 
    \[ \text{animal $\in$ ser vivo, cuadrado $\in$ figura }\]
        \item \emph{Programación Orientada a Objetos} (OOP)
    \end{itemize}
\end{frame}
%%%%%%%%%%%%%%%%%%%%%%%%%%%%%%%%%%%%%%%%%%%%%%%%%%%%%%%
\begin{frame}[fragile]\frametitle{Clases conocidas}
    \begin{minted}{python}
>>> type(pd.DataFrame()) # pandas
    <class 'pandas.core.frame.DataFrame'>
>>> type(np.linspace(0,1)) # numpy
    <class 'numpy.ndarray'>
>>> type(2) # base
    <class 'int'>
    \end{minted}
\end{frame}
%%%%%%%%%%%%%%%%%%%%%%%%%%%%%%%%%%%%%%%%%%%%%%%%%%%%%%%
\begin{frame}[fragile]\frametitle{Sintaxis}
    \begin{minted}{python}
class A:
   # Método constructor
   def __init__(self, n):
      self.n = n # son distintos!
      self.n2 = n**2
   def saludar(self): print("hola")

a = A(4) # Constructor
print(a.n, a.n2) # 4 16
print(type(a)) # <class '__main__.A'>
    \end{minted}

\end{frame}
%%%%%%%%%%%%%%%%%%%%%%%%%%%%%%%%%%%%%%%%%%%%%%%%%%%%%%%
\begin{frame}[fragile]\frametitle{Ejemplo: vector 2D}
    \begin{footnotesize}
    \begin{minted}{python}
class Vec2D:
   # Método constructor
   def __init__(self, x, y): # magic function
      self.x = x 
      self.y = y
   def __add__(self, v2): # magic function
      out = Vec2D(self.x, self.y) # copiar
      out.x += v2.x # agregar nuevos valores
      out.y += v2.y
      return out # retornar copia
   def __getitem__(self, i): # magic function
      if i == 0: return self.x
      if i == 1: return self.y
      else:
         print("ERROR") # i = 0 o 1
    \end{minted}
    \end{footnotesize}
\end{frame}
%%%%%%%%%%%%%%%%%%%%%%%%%%%%%%%%%%%%%%%%%%%%%%%%%%%%%%%
\darkSlide{Consola}
%%%%%%%%%%%%%%%%%%%%%%%%%%%%%%%%%%%%%%%%%%%%%%%%%%%%%%%
%%%%%%%%%%%%%%%%%%%%%%%%%%%%%%%%%%%%%%%%%%%%%%%%%%%%%%%
\section{Polimorfismo}
%%%%%%%%%%%%%%%%%%%%%%%%%%%%%%%%%%%%%%%%%%%%%%%%%%%%%%%
%%%%%%%%%%%%%%%%%%%%%%%%%%%%%%%%%%%%%%%%%%%%%%%%%%%%%%%
\begin{frame}\frametitle{Motivación}
    \begin{itemize}
        \item Abstraer "Figura"
    \[ \text{animal $\in$ ser vivo, cuadrado $\in$ figura }\]
        \item Sistema ecológico (pizarra)
    \end{itemize}
    $$ \text{Carnívoro}\quad\to\quad\text{Herbívoro}\quad\to\quad\text{Plantas} $$
\end{frame}
%%%%%%%%%%%%%%%%%%%%%%%%%%%%%%%%%%%%%%%%%%%%%%%%%%%%%%%
\begin{frame}[fragile]\frametitle{Sintaxis}
    \begin{minted}{python}
    class ClaseHeredada(ClaseBase):
        def funcion(self):
            # [...]
        # Opcional
        def __init__(self, params):
            super(ClaseHeredada, self).__init__(params)
    \end{minted}
\end{frame}
%%%%%%%%%%%%%%%%%%%%%%%%%%%%%%%%%%%%%%%%%%%%%%%%%%%%%%%
\begin{frame}[fragile]\frametitle{Ejemplo I (src: HowToGeek)}
    \begin{footnotesize}
    \begin{minted}{python}
class Bird:
  def intro(self): print("There are many types of birds.")
    
  def flight(self): print("Most of the birds can fly but some cannot.")
  
class sparrow(Bird):
  def flight(self): print("Sparrows can fly.")
    
class ostrich(Bird):
  def flight(self): print("Ostriches cannot fly.")
    
obj_bird = Bird(); obj_spr = sparrow(); obj_ost = ostrich()
obj_bird.flight(); obj_spr.flight(); obj_ost.flight()
    \end{minted}
    \end{footnotesize}
\end{frame}
%%%%%%%%%%%%%%%%%%%%%%%%%%%%%%%%%%%%%%%%%%%%%%%%%%%%%%%
\begin{frame}[fragile]\frametitle{Ejemplo II (src: HowToGeek)}
    \begin{footnotesize}
    \begin{minted}{python}
class Animal:
    def speak(self):
        raise NotImplementedError("Subclass must implement this method")

class Dog(Animal):
    def speak(self): return "Woof!"

class Cat(Animal):
    def speak(self): return "Meow!"

animals = [Dog(), Cat()]

for animal in animals:
    print(animal.speak())
    \end{minted}
    \end{footnotesize}
\end{frame}
%%%%%%%%%%%%%%%%%%%%%%%%%%%%%%%%%%%%%%%%%%%%%%%%%%%%%%%
%%%%%%%%%%%%%%%%%%%%%%%%%%%%%%%%%%%%%%%%%%%%%%%%%%%%%%%
\section{Documentación}
%%%%%%%%%%%%%%%%%%%%%%%%%%%%%%%%%%%%%%%%%%%%%%%%%%%%%%%
%%%%%%%%%%%%%%%%%%%%%%%%%%%%%%%%%%%%%%%%%%%%%%%%%%%%%%%
\begin{frame}\frametitle{Motivación}
    \begin{itemize}
        \item El código \emph{envejece}
        \item Uno \textbf{no} recuerda por qué hizo cosas
        \item Los futuros usuarios de nuestro código \huge{menos}
    \end{itemize}
\end{frame}
%%%%%%%%%%%%%%%%%%%%%%%%%%%%%%%%%%%%%%%%%%%%%%%%%%%%%%%
\begin{frame}[fragile]\frametitle{Docs de una función}
    \begin{minted}{python}
    def suma(x,y):
        """
        Esta función entrega la suma de sus argumentos. 
        args: 
            x: Un elemento sumable
            y: Un segundo elemento sumable
        Return: suma de entradas (Int)
        """
        z = x+y
        return z
    \end{minted}
\end{frame}
%%%%%%%%%%%%%%%%%%%%%%%%%%%%%%%%%%%%%%%%%%%%%%%%%%%%%%%
\begin{frame}[fragile]\frametitle{Docs de una función}
    \begin{minted}{python}
    >>> help(suma)
    \end{minted}

    \vspace{1cm}
    Se guarda en la variable mágica \code{suma.\_\_doc\_\_}

    \idea{Consola}
\end{frame}

%%%%%%%%%%%%%%%%%%%%%%%%%%%%%%%%%%%%%%%%%%%%%%%%%%%%%%%
\begin{frame}[fragile]\frametitle{Docs de una clase}
    \begin{footnotesize}
    \begin{minted}{python}
    class Vec2D:
        """
        Clase que representa un vector en R2.
        """
        def __init__(self,x,y):
            """
            Constructor, toma como parámetros coordenadas x e y.
            """
            self.x = x
            self.y = y
        def norma(self):
            """
            Entrega la norma euclideana del vector. Requiere 'sqrt'
            """
            return sqrt(self.x**2 + self.y**2)
    \end{minted}

    \vspace{-1cm}
    \hfill    \idea{Consola}
    \end{footnotesize}
\end{frame}
%%%%%%%%%%%%%%%%%%%%%%%%%%%%%%%%%%%%%%%%%%%%%%%%%%%%%%%
\begin{frame}[fragile]\frametitle{Documentación avanzada}
    Existe mucho software (Sphinx, Doxygen, etc). Veremos \code{Mkdocs}.

    \begin{minted}{bash}
    # pip install mkdocs
    # mkdocs new proyecto
    # cd proyecto
    # mkdocs serve
    # mkdocs build
    \end{minted} 

    \vspace{1cm}
    Se usa formato Markdown (popular por ser simple)

\end{frame}
%%%%%%%%%%%%%%%%%%%%%%%%%%%%%%%%%%%%%%%%%%%%%%%%%%%%%%%
\begin{frame}{Mkdocs}
    Veamos su documentación: \\https://www.mkdocs.org/getting-started/
\end{frame}
%%%%%%%%%%%%%%%%%%%%%%%%%%%%%%%%%%%%%%%%%%%%%%%%%%%%%%%
\begin{frame}[fragile]\frametitle{Otra opción}
    \begin{minted}{bash}
    # pip install pdoc3
    # cd proyecto
    # pdoc3 --html .
    \end{minted}

    \idea{Consola}
\end{frame}
%%%%%%%%%%%%%%%%%%%%%%%%%%%%%%%%%%%%%%%%%%%%%%%%%%%%%%%
\section{Testing}
%%%%%%%%%%%%%%%%%%%%%%%%%%%%%%%%%%%%%%%%%%%%%%%%%%%%%%%
\begin{frame}{Motivación}
    \begin{enumerate}
        \item Agrego cambio en función
        \item Pruebo nuevo código
        \item Falla un caso de borde escondidísimo
    \end{enumerate}

    \idea{Tener una cadena de pruebas automáticas}
\end{frame}
%%%%%%%%%%%%%%%%%%%%%%%%%%%%%%%%%%%%%%%%%%%%%%%%%%%%%%%
\begin{frame}{Solución: Pytest}
    \begin{center}
        \includegraphics[width=0.5\textwidth]{../images/pytest.png}
    \end{center}
\end{frame}
%%%%%%%%%%%%%%%%%%%%%%%%%%%%%%%%%%%%%%%%%%%%%%%%%%%%%%%
\begin{frame}[fragile]

    \begin{minted}{python}
    > myLib.py
    def suma(a,b): return a+b

    > test.py
    from myLib import suma
    
    def test_suma(): #debe comenzar con 'test'
        assert suma(1,2) == 3
    \end{minted}
    \begin{minted}{bash}
    $ pytest test.py
    \end{minted}
\end{frame}
%%%%%%%%%%%%%%%%%%%%%%%%%%%%%%%%%%%%%%%%%%%%%%%%%%%%%%%
\begin{frame}\frametitle{Recap}
    \begin{itemize}
        \item Documentación es fundamental
        \item Existen muchas herramientas para hacerlo en modo semi-automático
        \item \code{mkdocs} bonito pero trabajoso
        \item \code{pdoc3} feo pero fácil
        \item \code{git} es LA herramienta de desarrollo colaborativo
        \item Se basa en un grafo de estados
        \item Modificaciones locales: \code{add}, \code{commit}, \code{reset}
        \item Modificaciones remotas: \code{clone}, \code{pull}, \code{push}
    \end{itemize}
\end{frame}
%%%%%%%%%%%%%%%%%%%%%%%%%%%%%%%%%%%%%%%%%%%%%%%%%%%%%%%
\begin{frame}[noframenumbering]\frametitle{Mini ejercicios}
    \begin{enumerate}
        \item Instalar git
        \item Crear cuenta en \code{github.com}
        \item Subir llave SSH
        \item Crear repositorio en su cuenta
        \item Clonar el repositorio localmente
        \item Crear una carpeta de nombre 'proyecto'
        \item Dentro esa carpeta crear un archivo con dos funciones
        \item Instalar pdoc3 y correr \code{pdoc3 --html .} desde la carpeta
        \item Ver documentación en archivo html recién creado

        \item Hacer \emph{push} del cambio y verificar el cambio
    \end{enumerate}
    
    \begin{center}
        Turorial online: \code{https://www.w3schools.com/git/}
    \end{center}
\end{frame}
%%%%%%%%%%%%%%%%%%%%%%%%%%%%%%%%%%%%%%%%%%%%%%%%%%%%%%%
\begin{frame}\frametitle{Recap}
    \begin{itemize}
        \item Clases permiten crear abstracciones
        \item Todo objeto es una instancia de clase en Python
        \item Funciones mágicas permiten código legible!
    \end{itemize}
\end{frame}
%%%%%%%%%%%%%%%%%%%%%%%%%%%%%%%%%%%%%%%%%%%%%%%%%%%%%%%
\begin{frame}
    \maketitle
\end{frame}
%%%%%%%%%%%%%%%%%%%%%%%%%%%%%%%%%%%%%%%%%%%%%%%%%%%%%%%
\begin{frame}[fragile,noframenumbering]\frametitle{Mini ejercicios}
    \begin{footnotesize}
    \begin{itemize}
        \item Crear una clase Cuadrado que reciba como entrada el largo de sus aristas. Cree una función de clase que entregue el área del cuadrado.
        \item Extienda el ejemplo anterior al caso de un cubo.
        \item Instalar pdoc3 y correr \code{pdoc3 --html .} desde la carpeta y revisar documentación.
    \begin{minted}{python}
class Cuadrado:
   def __init__(self, lado): # [... implementación ...]
   def area(self):  # [... implementación ...]
    \end{minted}
        \item Explore el siguiente código:
    \begin{minted}{python}
class A:
    def a(self): print("Aa")

class B:
    def a(self):
        A.a(self)
        print("Ba")

b = B()
b.a()
        \end{minted}
    \end{itemize}
    \end{footnotesize}
\end{frame}
%%%%%%%%%%%%%%%%%%%%%%%%%%%%%%%%%%%%%%%%%%%%%%%%%%%%%%%

\end{document}
